% ===========================================================================
% Gói lệnh, định dạng và macro dùng chung cho toàn bộ đề án
% Biên dịch bằng XeLaTeX (cần cho tiếng Việt và phông Times New Roman)
% ===========================================================================

% --- Cỡ chữ cơ sở ---------------------------------------------------------
% Quy định trình bày luận văn yêu cầu cỡ chữ 13pt, nhưng lớp extreport không có
% tùy chọn này (xem ghi chú trong main.tex). Gói ``fontsize'' cho phép đặt cỡ
% chữ cơ sở tùy ý và tự tính lại toàn bộ các cỡ dẫn xuất (\large, \small,
% \footnotesize...) cùng khoảng cách dòng tương ứng.
\usepackage[fontsize=13pt]{fontsize}

% --- Ngôn ngữ và phông chữ ------------------------------------------------
\usepackage{fontspec}
% Ngôn ngữ chính là tiếng Việt (khai báo sau cùng), nạp kèm tiếng Anh cho các
% thuật ngữ và tài liệu tham khảo.
%
% Ghi chú: định dạng XeLaTeX trên máy biên dịch không nạp sẵn bộ quy tắc ngắt âm
% tiết, nên toàn bộ tài liệu được dựng ở chế độ không ngắt âm tiết. Với tiếng
% Việt đây là hành vi đúng (tiếng Việt không ngắt âm tiết trong từ). Để bù lại
% khả năng xuống dòng bị hạn chế, tham số \emergencystretch được nới rộng ở phần
% dưới, và danh mục tài liệu tham khảo được căn lề trái.
\usepackage[english,vietnamese]{babel}
\setmainfont{Times New Roman}[
  Ligatures=TeX,
  BoldFont={Times New Roman Bold},
  ItalicFont={Times New Roman Italic},
  BoldItalicFont={Times New Roman Bold Italic},
]
\setsansfont{Arial}[Ligatures=TeX]
\setmonofont{Consolas}[Scale=0.86]


% --- Trang và khổ giấy ----------------------------------------------------
% Lề theo quy định trình bày luận văn: trái 3,5cm (đóng quyển), phải 2cm
% headheight phải đủ chứa dòng đầu trang ở cỡ chữ 13pt; nếu thiếu, fancyhdr sẽ
% cảnh báo và dòng đầu trang bị đẩy chồng lên vùng chữ.
\usepackage[
  a4paper, top=2.5cm, bottom=2.5cm, left=3.5cm, right=2.0cm,
  headheight=18pt, headsep=0.6cm, footskip=1.2cm
]{geometry}
\usepackage{setspace}
\onehalfspacing                       % giãn dòng 1,5
\setlength{\parindent}{1cm}
\setlength{\parskip}{0.2\baselineskip}

% --- Toán học -------------------------------------------------------------
% Thứ tự nạp gói ở đây là bắt buộc: amsmath và mathtools phải được nạp TRƯỚC
% unicode-math, nếu không unicode-math sẽ báo lỗi định nghĩa lại ký hiệu.
\usepackage{amsmath}
\usepackage{mathtools}

% Phông toán tương thích Times.
% Mặc định LaTeX dựng công thức bằng Computer Modern, trong khi phần chữ dùng
% Times New Roman. Hai phông khác nhau rõ về nét và tỷ lệ nên công thức trông
% "lạc" giữa văn bản; ngoài ra Computer Modern chỉ có sẵn ở một số cỡ thiết kế
% cố định nên ở cỡ chữ 13pt LaTeX phải thay bằng cỡ gần nhất.
%
% TeX Gyre Termes Math là bộ ký hiệu toán OpenType dựng theo Times, dùng trực
% tiếp với XeLaTeX qua gói unicode-math. Gói này đã bao gồm toàn bộ ký hiệu của
% amssymb/amsfonts nên KHÔNG nạp thêm hai gói đó (sẽ xung đột định nghĩa).
% Tùy chọn ``math-style=ISO'' đặt chữ cái Latin và Hy Lạp trong công thức ở dạng
% nghiêng theo chuẩn ISO 80000-2.
\usepackage{unicode-math}
\setmathfont{TeX Gyre Termes Math}
\unimathsetup{math-style=ISO, bold-style=ISO}

% unicode-math cung cấp \symbf thay cho gói bm; giữ lại tên \bm cho tương thích
% với phần nội dung đã viết.
\newcommand{\bm}[1]{\symbf{#1}}

\numberwithin{equation}{chapter}

% --- Bảng biểu ------------------------------------------------------------
\usepackage{booktabs}
\usepackage{longtable}
\usepackage{tabularx}
\usepackage{multirow}
\usepackage{array}
\usepackage{makecell}
\usepackage{colortbl}
\renewcommand{\arraystretch}{1.22}
\newcolumntype{L}[1]{>{\raggedright\arraybackslash}p{#1}}
\newcolumntype{C}[1]{>{\centering\arraybackslash}p{#1}}
\newcolumntype{R}[1]{>{\raggedleft\arraybackslash}p{#1}}

% --- Hình ảnh và đồ họa ---------------------------------------------------
\usepackage{graphicx}
\graphicspath{{figures/}{../results/figures/}}
\usepackage{float}
\usepackage{subcaption}
\usepackage{xcolor}
\usepackage{tikz}
\usetikzlibrary{
  shapes.geometric, shapes.misc, arrows.meta, positioning, calc,
  fit, backgrounds, decorations.pathreplacing, patterns, chains, shadows.blur
}
\usepackage{pgfplots}
\pgfplotsset{compat=1.18}

% --------------------------------------------------------------------------
% Tự động thu nhỏ hình/sơ đồ vượt khổ chữ.
%
% Môi trường ``fitfigure'' đo chiều rộng thực tế của nội dung rồi CHỈ thu nhỏ
% khi nội dung rộng hơn khổ chữ. Đây là điểm khác biệt quan trọng so với
% \resizebox{\textwidth}{!}{...}: lệnh đó ép MỌI nội dung về đúng chiều rộng khổ
% chữ, khiến những sơ đồ vốn đã vừa bị phóng to quá mức và cỡ chữ trong hình trở
% nên không đồng nhất giữa các hình.
%
% Hiện thực bằng hộp lưu (savebox) của LaTeX, không cần thêm gói phụ thuộc.
% --------------------------------------------------------------------------
\newsavebox{\cffitbox}
\newenvironment{fitfigure}
  {\begin{lrbox}{\cffitbox}}
  {\end{lrbox}%
   % Ghi chiều rộng tự nhiên của hình ra tệp nhật ký. Nhờ đó có thể rà soát
   % những sơ đồ rộng hơn khổ chữ nhiều lần - đó là các sơ đồ sẽ bị thu nhỏ
   % mạnh và chữ bên trong trở nên khó đọc, cần thiết kế lại cho gọn hơn.
   \typeout{FITFIG \the\wd\cffitbox\space / \the\textwidth}%
   \ifdim\wd\cffitbox>\textwidth
     \resizebox{\textwidth}{!}{\usebox{\cffitbox}}%
   \else
     \usebox{\cffitbox}%
   \fi}

% --- Chú thích hình/bảng --------------------------------------------------
\usepackage[font=small,labelfont=bf,justification=centering,skip=6pt]{caption}

% --- Danh sách ------------------------------------------------------------
\usepackage{enumitem}
\setlist{itemsep=2pt, topsep=4pt, parsep=0pt}

% --- Mã nguồn -------------------------------------------------------------
\usepackage{listings}
\definecolor{codebg}{HTML}{F7F8FA}
\definecolor{codekw}{HTML}{1F4E79}
\definecolor{codestr}{HTML}{548235}
\definecolor{codecmt}{HTML}{8A8A8A}
\definecolor{codenum}{HTML}{A0A0A0}
\lstdefinestyle{cf}{
  backgroundcolor=\color{codebg},
  basicstyle=\ttfamily\scriptsize,
  keywordstyle=\color{codekw}\bfseries,
  stringstyle=\color{codestr},
  commentstyle=\color{codecmt}\itshape,
  numberstyle=\tiny\color{codenum},
  numbers=left, numbersep=7pt,
  breaklines=true, breakatwhitespace=false,
  showstringspaces=false, captionpos=b,
  frame=single, rulecolor=\color{gray!35},
  xleftmargin=14pt, framexleftmargin=14pt,
  tabsize=4, columns=fullflexible, keepspaces=true,
  literate=%
    {á}{{\'a}}1 {à}{{\`a}}1 {ả}{{\h a}}1 {ã}{{\~a}}1 {ạ}{{\d a}}1
    {ă}{{\u a}}1 {â}{{\^a}}1 {đ}{{\dj}}1 {é}{{\'e}}1 {è}{{\`e}}1
    {ê}{{\^e}}1 {í}{{\'i}}1 {ì}{{\`i}}1 {ó}{{\'o}}1 {ò}{{\`o}}1
    {ô}{{\^o}}1 {ơ}{{\o'}}1 {ú}{{\'u}}1 {ù}{{\`u}}1 {ư}{{\u'}}1
    {ý}{{\'y}}1
}
\lstset{style=cf}
\renewcommand{\lstlistingname}{Đoạn mã}

% --- Tham chiếu chéo và siêu liên kết -------------------------------------
\usepackage[
  colorlinks=true, linkcolor=black, citecolor=black,
  urlcolor=blue!55!black, filecolor=black, bookmarksnumbered=true
]{hyperref}
\usepackage{bookmark}
\usepackage{url}
\urlstyle{same}

% --- Tiêu đề chương/mục ---------------------------------------------------
\usepackage{titlesec}
\titleformat{\chapter}[display]
  {\normalfont\LARGE\bfseries\centering}
  {\chaptertitlename\ \thechapter}{14pt}{\LARGE}
\titlespacing*{\chapter}{0pt}{-18pt}{26pt}
\titleformat{\section}{\normalfont\large\bfseries}{\thesection}{0.8em}{}
\titleformat{\subsection}{\normalfont\normalsize\bfseries}{\thesubsection}{0.8em}{}
\titleformat{\subsubsection}{\normalfont\normalsize\itshape}{\thesubsubsection}{0.8em}{}
\setcounter{secnumdepth}{3}
\setcounter{tocdepth}{2}

% --- Đầu trang / chân trang -----------------------------------------------
\usepackage{fancyhdr}
\pagestyle{fancy}
\fancyhf{}
\fancyhead[R]{\small\nouppercase{\leftmark}}
\fancyfoot[C]{\thepage}
\renewcommand{\headrulewidth}{0.4pt}
\fancypagestyle{plain}{\fancyhf{}\fancyfoot[C]{\thepage}%
  \renewcommand{\headrulewidth}{0pt}}

% --- Danh mục viết tắt ----------------------------------------------------
\usepackage{longtable}

% ===========================================================================
% MÀU VÀ KIỂU TIKZ DÙNG CHUNG CHO CÁC SƠ ĐỒ KIẾN TRÚC
% ===========================================================================
\definecolor{cfBlue}{HTML}{1F4E79}
\definecolor{cfBlueL}{HTML}{DCE9F5}
\definecolor{cfOrange}{HTML}{C55A11}
\definecolor{cfOrangeL}{HTML}{FDF0E0}
\definecolor{cfGreen}{HTML}{548235}
\definecolor{cfGreenL}{HTML}{E8F0E0}
\definecolor{cfRed}{HTML}{C0504D}
\definecolor{cfRedL}{HTML}{F7E4E3}
\definecolor{cfGray}{HTML}{5A5A5A}
\definecolor{cfGrayL}{HTML}{F0F0F0}
\definecolor{cfPurple}{HTML}{7030A0}
\definecolor{cfPurpleL}{HTML}{EDE3F5}

\tikzset{
  cfnode/.style={
    draw, rounded corners=3pt, align=center, font=\small,
    minimum height=9mm, minimum width=24mm, inner sep=4pt, line width=0.7pt
  },
  source/.style={cfnode, fill=cfBlueL, draw=cfBlue},
  stream/.style={cfnode, fill=cfOrangeL, draw=cfOrange},
  compute/.style={cfnode, fill=cfGreenL, draw=cfGreen},
  storage/.style={cfnode, fill=cfRedL, draw=cfRed},
  external/.style={cfnode, fill=cfPurpleL, draw=cfPurple},
  ui/.style={cfnode, fill=cfGrayL, draw=cfGray},
  decision/.style={
    draw=cfBlue, fill=cfBlueL, diamond, aspect=2.2, align=center,
    font=\scriptsize, inner sep=2pt, line width=0.7pt
  },
  flowarrow/.style={-{Stealth[length=2.4mm]}, line width=0.75pt, draw=cfGray},
  dataflow/.style={
    -{Stealth[length=2.4mm]}, line width=0.75pt, draw=cfBlue,
    font=\scriptsize
  },
  faultflow/.style={
    -{Stealth[length=2.4mm]}, line width=0.75pt, draw=cfRed, dashed,
    font=\scriptsize
  },
  grouplabel/.style={font=\scriptsize\bfseries, text=cfGray},
  zone/.style={
    draw=gray!45, dashed, rounded corners=5pt, inner sep=7pt, line width=0.6pt
  },
}

% ===========================================================================
% MACRO TIỆN ÍCH
% ===========================================================================
% Định dạng số theo chuẩn Việt Nam (dấu phẩy thập phân)
\usepackage{siunitx}
\sisetup{
  output-decimal-marker={,},
  group-separator={.},
  group-minimum-digits=4,
  detect-all
}

% --------------------------------------------------------------------------
% Chèn hình an toàn: nếu hình chưa được sinh (chưa chạy hết thực nghiệm),
% thay bằng một khung báo thiếu thay vì làm hỏng toàn bộ quá trình biên dịch.
% --------------------------------------------------------------------------
% Tên file chứa dấu gạch dưới nên phải dùng \detokenize để in nguyên văn,
% tránh việc LaTeX hiểu "_" là ký hiệu chỉ số dưới trong chế độ toán.
% Hình được đặt trong \makebox rộng đúng bằng khổ chữ và bắt đầu bằng \noindent:
% nếu không, phần thụt đầu dòng 1cm sẽ cộng vào chiều rộng hình và làm tràn lề
% ngay cả khi hình được khai báo hẹp hơn khổ chữ.
\newcommand{\cffig}[2][\textwidth]{%
  \par\noindent
  \makebox[\linewidth][c]{%
    \IfFileExists{../results/figures/#2.pdf}
      {\includegraphics[width=#1]{#2}}
      {\fbox{\parbox[c][2.8cm][c]{0.8\textwidth}{\centering\small\itshape
        Hình \texttt{\detokenize{#2}} chưa được sinh --
        chạy lại các script thực nghiệm để tạo hình này.}}}%
  }\par
}
% Tên gọi ngắn dùng trong các chương
\providecommand{\fig}{}
\renewcommand{\fig}[2][\textwidth]{\cffig[#1]{#2}}

% Thuật ngữ tiếng Anh in nghiêng khi xuất hiện lần đầu
\newcommand{\eng}[1]{\textit{#1}}
% Tên thành phần hệ thống
\newcommand{\comp}[1]{\textsf{#1}}

% --------------------------------------------------------------------------
% Tên biến / đường dẫn mã nguồn.
%
% Chuỗi kiểu ``scripts/06_verify_consistency.py'' là một khối liền không có
% khoảng trắng nên TeX không thể ngắt dòng, gây tràn lề khi chuỗi rơi vào cuối
% dòng. Định nghĩa dưới đây tạo điểm ngắt hợp lệ ngay sau mỗi dấu gạch dưới
% bằng cách định nghĩa lại ``\_'' trong phạm vi cục bộ của lệnh.
% --------------------------------------------------------------------------
\newcommand{\code}[1]{{%
  \ttfamily\small
  \def\_{\char`\_\allowbreak}%
  #1}}

% --------------------------------------------------------------------------
% Cho phép TeX co giãn khoảng trắng nhiều hơn trước khi chấp nhận tràn lề.
% Cần thiết với văn bản tiếng Việt xen lẫn thuật ngữ tiếng Anh và tên riêng
% nước ngoài: những từ này không có quy tắc ngắt âm tiết tiếng Việt nên trình
% bày dễ bị tràn ở các dòng có từ dài.
% --------------------------------------------------------------------------
\emergencystretch=3em
\hbadness=2500        % bỏ qua cảnh báo dòng thưa ở mức chấp nhận được
\tolerance=1200

% --------------------------------------------------------------------------
% Tham số đặt hình/bảng. Mặc định của LaTeX quá dè dặt: một hình cỡ trung bình
% bị đẩy sang "trang chỉ có hình", để lại hơn nửa trang trắng. Nới rộng tỷ lệ
% hình được phép chiếm ở đầu/cuối trang và siết \floatpagefraction để hình chỉ
% chiếm trọn một trang khi thực sự đủ lớn.
% --------------------------------------------------------------------------
\renewcommand{\topfraction}{0.90}
\renewcommand{\bottomfraction}{0.70}
\renewcommand{\textfraction}{0.08}
\renewcommand{\floatpagefraction}{0.80}
\setcounter{topnumber}{3}
\setcounter{bottomnumber}{2}
\setcounter{totalnumber}{4}

% Ô nhấn mạnh kết quả chính
\usepackage[most]{tcolorbox}
\newtcolorbox{keyresult}[1][]{
  enhanced, breakable,
  colback=cfBlueL!45, colframe=cfBlue, boxrule=0.7pt, arc=3pt,
  left=8pt, right=8pt, top=6pt, bottom=6pt,
  fonttitle=\bfseries\small, title={Kết quả chính}, #1
}
\newtcolorbox{notebox}[1][]{
  enhanced, breakable,
  colback=cfOrangeL!60, colframe=cfOrange, boxrule=0.7pt, arc=3pt,
  left=8pt, right=8pt, top=6pt, bottom=6pt,
  fonttitle=\bfseries\small, title={Lưu ý}, #1
}

% --------------------------------------------------------------------------
% Trích dẫn tài liệu tham khảo
% Danh mục tài liệu được đánh số thủ công bằng môi trường enumerate (phù hợp
% quy định trình bày của cơ sở đào tạo). Lệnh \cite được định nghĩa lại để tham
% chiếu tới nhãn tương ứng, nhờ đó số thứ tự trích dẫn luôn khớp với danh mục
% ngay cả khi bổ sung hoặc sắp xếp lại tài liệu.
% --------------------------------------------------------------------------
% Hỗ trợ trích dẫn nhiều tài liệu cùng lúc: \cite{a,b,c} -> [8, 9, 10]
\makeatletter
\newcommand{\cf@citeone}[1]{\ref{ref:#1}}
\newcommand{\cf@citelist}[1]{%
  \def\cf@sep{}%
  \@for\cf@key:=#1\do{\cf@sep\cf@citeone{\cf@key}\def\cf@sep{, }}%
}
\renewcommand{\cite}[1]{[\cf@citelist{#1}]}
\makeatother

% Ký hiệu toán học hay dùng
\newcommand{\PD}{\mathrm{PD}}
\newcommand{\LGD}{\mathrm{LGD}}
\newcommand{\EAD}{\mathrm{EAD}}
\newcommand{\EL}{\mathrm{EL}}
\newcommand{\WOE}{\mathrm{WOE}}
\newcommand{\IV}{\mathrm{IV}}
\newcommand{\PSI}{\mathrm{PSI}}
\newcommand{\AUC}{\mathrm{AUC}}
\newcommand{\KS}{\mathrm{KS}}
\newcommand{\DTI}{\mathrm{DTI}}
\newcommand{\LTV}{\mathrm{LTV}}
\DeclareMathOperator*{\argmin}{arg\,min}
\DeclareMathOperator*{\argmax}{arg\,max}
\DeclareMathOperator{\sigmoid}{\sigma}
\DeclareMathOperator{\Ex}{\mathbb{E}}

% Tránh chữ mồ côi
\widowpenalty=10000
\clubpenalty=10000
\raggedbottom
